\nocite{Chiesa2024-fopp-course}
\begin{definition}
\label{QESAT}
\lean{QESAT}
\leanok
Let \(m,n >0\) be integers, \(\F\) be a field, \(R = \F[x_{1},\dots,x_{n}]\) be its polynomial ring and \(R_{d}\) denote the subspace of polynomials of degree \(\leq d\).

\[
  \QESAT(\F,n) :=
  \{ (p_1, \ldots, p_m) \in R^m_{2} \mid
    \exists a \in \F^n, \,
    \forall i \in [m], \,
    p_i(a_1, \ldots, a_n) = 0 \}.
\]
\end{definition}

For example, $(x + 1, xy + z) \in \QESAT(\F_2,3)$.
The goal of this chapter is to construct an exponential-length PCP for
\(\QESAT(\F_2,n)\) (Theorem \ref{QESAT_exp_PCP}). Since \(\QESAT(\F_2,n)\) is NP-complete, this will imply the existence of exponential-length PCPs for all languages in NP.

\section{Oracle computations and verifiers}
Instead of oracle Turing machines, we will define verifiers using Lean monads, specifically the \texttt{OracleComp} free monad from VCV-io.

\begin{definition}[Oracle specification]
\label{OracleSpec}
\lean{OracleSpec}
\leanok
An \emph{oracle specification} defines the type of the queries and outputs that an oracle can handle.
\end{definition}

\begin{definition}
\label{randOracleSpec}
\lean{randOracleSpec}
\uses{OracleSpec}
\leanok
For this section, the \emph{randomness oracle specification} will be of a monadic function that takes
nothing and returns an element from a field \(\F\).
\end{definition}

\begin{definition}[Query implementation]
\label{QueryImpl}
\lean{QueryImpl}
\leanok
A \emph{query implementation} is a monadic function whose behavior conforms to a given oracle specification.
\end{definition}

\begin{definition}
\label{uniformSample}
\uses{QueryImpl, randOracleSpec}
\lean{uniformSample}
\leanok
The \emph{uniform sampling query implementation} for a field \(\F\) is a (mathematical) 
query implementation for the randomness oracle specification that samples uniformly from \(\F\).
\end{definition}
This is the the oracle implementation that will be used to prove theorems about our verifiers,
since it is compatible with the more classical notion of probability distributions used in the previous chapters.
\begin{remark}[Implementing verifiers]
If we want to actually run the verifiers, we can also use concrete implementations based on the IO monad
and system randomness for pseudorandom sampling.
\end{remark}
\begin{definition}[Oracle context]
\label{OracleContext}
\lean{OracleContext}
\uses{OracleSpec, QueryImpl}
\leanok
An \emph{oracle context} is a bundle consisting of an oracle specification and a query implementation for that specification.
\end{definition}
\begin{definition}[Randomness oracle]
\label{randOracle}
\uses{OracleContext, randOracleSpec, uniformSample}
\lean{randOracle} \leanok
As \emph{randomness oracle}, we will use the {uniform sampling query implementation} for the randomness oracle specification.
\end{definition}
\begin{definition}[Oracle computation]
\label{OracleComp}
\uses{OracleSpec}
\lean{OracleComp}\leanok
An \emph{oracle computation} is a monadic computation in the free monad generated by an oracle specification,
i.e. a description (a ``script'') of a program consisting of a chain of oracle queries (compatible with any query implementation of the specification)
and deterministic Lean computations.
\end{definition}

A \emph{verifier} is a function that takes an input \(x \in \{0,1\}^*\) and produces an oracle computation
(with query access to a randomness oracle) that outputs a boolean. This definition can be extended to include
other oracles (typically a proof oracle). In Lean this convention is represented directly by function types returning
\texttt{OracleComp} computations, rather than by a standalone verifier declaration.

We also need to define verifiers that do not take any inputs (for example the BLR verifier).
Since they use the same randomness oracle specification as the normal verifiers, they can
compose nicely with them.
A \emph{unit verifier} (or just \emph{verifier} when it is clear from the context) is an oracle computation with query access to a randomness oracle that outputs a boolean.
This definition can be extended to include
other oracles (typically a proof oracle).

\begin{remark}[Time complexity]
The price to pay for using these definitions for verifiers instead of Turing machines is that we have no simple way to enforce
the time complexity of the verifier. For that, a possible approach would be to restrict the allowed operations to a strict set
of operations instead of general Lean computations, and track the time cost using query complexity.
\end{remark}
\section{PCP and LPCP definitions}

\begin{definition}
\label{PCPVerifier}
\lean{PCPVerifier}
\uses{OracleComp,randOracleSpec}
\leanok
A \emph{PCP verifier} is a verifier with query access to a function $\pi : [\ell (|x|)] \to \Sigma$.
\end{definition}

\begin{definition}
\uses{PCPVerifier}
\label{PCP}
\lean{PCP}
\leanok
A language $L \subseteq \{0,1\}^*$ is in
$\PCP[\varepsilon_c, \varepsilon_s, \Sigma, \ell, q, r]$
if there exists a polynomial-time PCP verifier $V$
such that for every $x \in \{0,1\}^*$,
$V$ makes at most $q(|x|)$ queries to the proof oracle and
at most $r(|x|)$ randomness queries, and the following holds:
\begin{itemize}
  \item Completeness: If $x \in L$, then
  $\exists \pi \in \Sigma^{\ell(|x|)}, \,
  \Pr\left[V^\pi(x)=1 \right] \geq 1 - \varepsilon_c$.
  \item Soundness: If $x \notin L$, then
  $\forall \widetilde{\pi} \in \Sigma^{\ell(|x|)},\,
  \Pr\left[V^{\widetilde{\pi}}(x)=1 \right] \leq \varepsilon_s$.
\end{itemize}
\end{definition}

\begin{definition}
\label{LPCPVerifier}
\lean{LPCPVerifier}
\uses{OracleComp,randOracleSpec}
\leanok
A \emph{LPCP verifier} is a verifier with query access to a linear function
$\langle \pi,\cdot \rangle : \Sigma^{\ell(|x|)} \to \Sigma$.
\end{definition}

\begin{definition}
\uses{LPCPVerifier}
\label{LPCP}
\lean{LPCP}
\leanok
A language $L \subseteq \{0,1\}^*$ is in
$\LPCP[\varepsilon_c, \varepsilon_s, \Sigma, \ell, q, r]$
if there exists a polynomial-time LPCP verifier $V$
such that for every $x \in \{0,1\}^*$,
$V$ makes at most $q(|x|)$ queries to the linear proof oracle and
at most $r(|x|)$ randomness queries, and the following holds:
\begin{itemize}
  \item Completeness: If $x \in L$, then
  $\exists \pi \in \Sigma^{\ell(|x|)}, \,
  \Pr\left[V^{\langle \pi, \cdot \rangle}(x)=1 \right]
  \geq 1 - \varepsilon_c$.
  \item Soundness: If $x \notin L$, then
  $\forall \widetilde{\pi} \in \Sigma^{\ell(|x|)},\,
  \Pr\left[V^{\langle \widetilde{\pi}, \cdot \rangle}(x)=1 \right]
  \leq \varepsilon_s$.
\end{itemize}
\end{definition}
